\documentclass[../../main.tex]{subfiles}% 注意这里的文件路径不能用 ./main.tex ,否则用latexmk编译子文件会报错
\graphicspath{{\subfix{./image/}}} % 指定图片目录,后续可以直接使用图片文件名
% 注意这里的文件路径不能用 ../../image/ ,否则用latexmk编译子文件会报错

% 例如:
% \begin{figure}[H]
% \centering
% \includegraphics[scale=0.3]{图.png}
% \caption{}
% \label{figure:图}
% \end{figure}
% 注意:上述\label{}一定要放在\caption{}之后,否则引用图片序号会只会显示??.

\begin{document}

\section{Abel变换}

\begin{theorem}[Abel变换]\label{theorem:Abel变换}
设\(\{a_n\}_{n = 1}^{N}\),\(\{b_n\}_{n = 1}^{N}\)是数列,则有恒等式
\begin{align*}
\sum\limits_{k = 1}^{N} a_k b_k = \sum\limits_{k = 1}^{N - 1} (a_k - a_{k + 1}) \sum\limits_{i = 1}^{k} b_i + a_N \sum\limits_{k = 1}^{N} b_k.
\end{align*}
进而也有
\begin{align*}
\sum_{k=1}^N{a_k\left( b_k-b_{k+1} \right)}=a_1b_1-a_Nb_{N+1}+\sum_{k=1}^{N-1}{\left( a_{k+1}-a_k \right) b_{k+1}}.
\end{align*}
\end{theorem}
\begin{note}
\hyperref[theorem:Abel变换]{Abel变换}的证明想法“强行裂项”是一种很重要的思想.
\end{note}
\begin{proof}

为了计算\(\sum\limits_{j = 1}^{N - 1} (a_j - a_{j + 1}) \sum\limits_{i = 1}^{j} b_i + a_N \sum\limits_{i = 1}^{N} b_i\),我们来强行构造裂项,差什么就给他补上去再补回来,即:
\begin{align*}
&\sum\limits_{j = 1}^{N - 1} (a_j - a_{j + 1}) \sum\limits_{i = 1}^{j} b_i + a_N \sum\limits_{i = 1}^{N} b_i = \sum\limits_{j = 1}^{N - 1} \left(a_j \sum\limits_{i = 1}^{j} b_i - a_{j + 1} \sum\limits_{i = 1}^{j} b_i\right) + a_N \sum\limits_{i = 1}^{N} b_i
\\
&= \sum\limits_{j = 1}^{N - 1} \left(a_j \sum\limits_{i =1}^{j} b_i - a_{j + 1} \sum\limits_{i = 1}^{j + 1} b_i\right) + \sum\limits_{j = 1}^{N - 1} \left(a_{j + 1} \sum\limits_{i = 1}^{j + 1} b_i - a_{j + 1} \sum\limits_{i = 1}^{j} b_i\right) + a_N \sum\limits_{i = 1}^{N} b_i
\\
&= a_1 b_1 - a_N \sum\limits_{i = 1}^{N} b_i + \sum\limits_{j = 1}^{N - 1} a_{j + 1} b_{j + 1} + a_N \sum\limits_{i = 1}^{N} b_i
= \sum\limits_{j = 1}^{N} a_j b_j.
\end{align*}
再利用上式直接计算即得
\begin{align*}
\sum_{k=1}^N{a_k\left( b_k-b_{k+1} \right)}=a_1b_1-a_Nb_{N+1}+\sum_{k=1}^{N-1}{\left( a_{k+1}-a_k \right) b_{k+1}}.
\end{align*}

\end{proof}

\begin{proposition}[经典乘积极限结论]\label{proposition:经典乘积极限结论}
设\(a_1 \geqslant a_2 \geqslant \cdots \geqslant a_n \geqslant 0\)且\(\lim_{n \to \infty} a_n = 0\),极限\(\lim_{n \to \infty} \sum\limits_{k = 1}^{n} a_k b_k\)存在.证明
\begin{align*}
\lim_{n \to \infty} (b_1 + b_2 + \cdots + b_n)a_n = 0. 
\end{align*}
\end{proposition}
\begin{note}
为了估计\(\sum\limits_{j = 1}^{n} b_j\),前面的有限项不影响.而要用上极限\(\sum\limits_{n = 1}^{\infty} a_n b_n\)收敛,自然想到\(\sum\limits_{j = 1}^{n} b_j = \sum\limits_{j = 1}^{n} \frac{b_j a_j}{a_j}\)和\hyperref[theorem:Abel变换]{Abel变换}.而\(a_j\)的单调性能用在\hyperref[theorem:Abel变换]{Abel变换}之后去绝对值.
\end{note}
\begin{proof}

不妨设\(a_1 \geqslant a_2 \geqslant \cdots \geqslant a_n > 0\).则由于级数\(\sum\limits_{n = 1}^{\infty} a_n b_n\)收敛,故对$\forall \varepsilon >0,$存在\(N \in \mathbb{N}\),使得
\begin{align*}
\left|\sum\limits_{i = N + 1}^{m} a_i b_i\right| \leqslant \varepsilon, \forall m \geqslant N + 1.
\end{align*}
当\(n \geqslant N + 1\),由\hyperref[theorem:Abel变换]{Abel变换},我们有
\begin{align*}
&\left|\sum\limits_{j = N + 1}^{n} b_j\right| = \left|\sum\limits_{j = N + 1}^{n} \frac{a_j b_j}{a_j}\right|
= \left|\sum\limits_{j = N + 1}^{n - 1} \left(\frac{1}{a_j} - \frac{1}{a_{j + 1}}\right) \sum\limits_{i = N + 1}^{j} a_i b_i + \frac{1}{a_n} \sum\limits_{i = N + 1}^{n} a_i b_i\right|
\\
&\leqslant \sum\limits_{j = N + 1}^{n - 1} \left(\left|\frac{1}{a_j} - \frac{1}{a_{j + 1}}\right| \cdot \left|\sum\limits_{i = N + 1}^{j} a_i b_i\right|\right) + \frac{1}{|a_n|} \left|\sum\limits_{i = N + 1}^{n} a_i b_i\right|
\\
&\leqslant \left|\sum\limits_{i = N + 1}^{n} a_i b_i\right|\cdot\sum\limits_{j = N + 1}^{n - 1} \left(\left|\frac{1}{a_j} - \frac{1}{a_{j + 1}}\right| \right) + \frac{1}{|a_n|} \left|\sum\limits_{i = N + 1}^{n} a_i b_i\right|
\\
&\leqslant \varepsilon \left[\sum\limits_{j = N + 1}^{n - 1} \left(\frac{1}{a_{j + 1}} - \frac{1}{a_j}\right) + \frac{1}{a_n}\right]
= \varepsilon \left(\frac{2}{a_n} - \frac{1}{a_{N + 1}}\right).
\end{align*}
因此我们有
\begin{align*}
&\underset{n\rightarrow \infty}{{\varlimsup }} \left|a_n \sum\limits_{j = 1}^{n} b_j\right| \leqslant \underset{n\rightarrow \infty}{{\varlimsup }} \left|a_n \sum\limits_{j = 1}^{N} b_j\right| + \underset{n\rightarrow \infty}{{\varlimsup }} \left|a_n \sum\limits_{j = N + 1}^{n} b_j\right|
\leqslant \underset{n\rightarrow \infty}{{\varlimsup }} \left|a_n \sum\limits_{j = 1}^{N} b_j\right| + \varepsilon \underset{n\rightarrow \infty}{{\varlimsup }} \left(2 - \frac{a_n}{a_{N + 1}}\right)= 2\varepsilon.
\end{align*}
由\(\varepsilon\)任意性即可得
\(\underset{n\rightarrow \infty}{{\varlimsup }} \left|a_n \sum\limits_{j = 1}^{n} b_j\right| = 0\),
于是就证明了
\(\lim_{n \to \infty} (b_1 + b_2 + \cdots + b_n)a_n = 0\).

\end{proof}

\begin{example}
设有收敛级数$\sum\limits_{n=1}^{\infty} a_n = S$，证明：
\begin{enumerate}[(1)]
\item $\lim\limits_{n\to\infty} \dfrac{a_1 + 2a_2 + \cdots + na_n}{n} = 0$；
\item 级数$\sum\limits_{n=1}^{\infty} \dfrac{a_1 + 2a_2 + \cdots + na_n}{n(n+1)}$收敛，且其和也是$S$.
\end{enumerate}
\end{example}
\begin{note}
本题$a_n$不一定非负，不能用\hyperref[theorem:级数的Fubini定理]{级数的Fubini定理}直接对级数换序.
\end{note}
\begin{proof}
\begin{enumerate}[(1)]
\item 令$S_n = \sum\limits_{k=1}^n a_k, \forall n\in\mathbb{N}$。则由\hyperref[theorem:Stolz定理]{Stolz公式}可得
\begin{align*}
\lim_{n\to\infty} \frac{\sum\limits_{k=1}^{n-1} S_k}{n} = \lim_{n\to\infty} S_n = S.
\end{align*}
再由\hyperref[theorem:Abel变换]{Abel变换}可得
\begin{align*}
\frac{\sum\limits_{k=1}^n ka_k}{n} = \frac{n\sum\limits_{k=1}^n a_k - \sum\limits_{k=1}^{n-1} \sum\limits_{i=1}^k a_i}{n} = S_n - \frac{\sum\limits_{k=1}^{n-1} S_k}{n}, \quad \forall n\in\mathbb{N}.
\end{align*}
令$n\to\infty$得
\begin{align*}
\lim_{n\to\infty} \frac{\sum\limits_{k=1}^n ka_k}{n} = \lim_{n\to\infty} S_n - \lim_{n\to\infty} \frac{\sum\limits_{k=1}^{n-1} S_k}{n} = S - S = 0.
\end{align*}

\item 对$\forall N\in\mathbb{N}$，利用\hyperref[theorem:Abel变换]{Abel变换}得
\begin{align*}
\sum_{n=1}^N \frac{\sum\limits_{k=1}^n ka_k}{n(n+1)}
= \sum_{n=1}^N \left( \sum_{k=1}^n ka_k \right) \left( \frac{1}{n} - \frac{1}{n+1} \right) 
= a_1 - \frac{\sum\limits_{k=1}^N ka_k}{N+1} + \sum_{n=1}^{N-1} a_{n+1} 
= a_1 - \frac{\sum\limits_{k=1}^N ka_k}{N+1} + \sum_{n=2}^N a_n.
\end{align*}
由(1)知$\lim\limits_{N\to\infty} \frac{\sum\limits_{k=1}^N ka_k}{N+1} = 0$，故对上式令$N\to\infty$得
\begin{align*}
\lim_{N\to\infty} \sum_{n=1}^N \frac{\sum\limits_{k=1}^n ka_k}{n(n+1)}
= a_1 - \lim_{N\to\infty} \frac{\sum\limits_{k=1}^N ka_k}{N+1} + \lim_{N\to\infty} \sum_{n=2}^N a_n 
= a_1 + \lim_{N\to\infty} \sum_{n=1}^N a_n - a_1 = S.
\end{align*}
\end{enumerate}

\end{proof}

\begin{example}
设级数 $\sum\limits_{n=1}^{\infty} na_n$ 收敛，证明：
\begin{enumerate}[(1)]
\item 对任一正整数 $k$，级数 $\sum\limits_{n=1}^{\infty} na_{n+k}$ 收敛；
\item $\lim\limits_{k\to\infty} \sum\limits_{n=1}^{\infty} na_{n+k} = 0.$
\end{enumerate}
\end{example}
\begin{proof}
\begin{enumerate}[(1)]
\item 对$\forall k\in \mathbb{N}$，注意到
\begin{align*}
\sum\limits_{n=1}^{\infty} n a_{n+k} = \sum\limits_{n=k+1}^{\infty} (n-k) a_n = \sum\limits_{n=k+1}^{\infty} n a_n - k \sum\limits_{n=k+1}^{\infty} a_n = \sum\limits_{n=k+1}^{\infty} n a_n - k \sum\limits_{n=k+1}^{\infty} \frac{n a_n}{n}.
\end{align*}
由题设知$\sum\limits_{n=k+1}^{\infty} n a_n$收敛，又由Dirichlet判别法知$\sum\limits_{n=k+1}^{\infty} \frac{n a_n}{n}$收敛。故$\sum\limits_{n=1}^{\infty} n a_{n+k}$收敛。

\item 对$\forall k\in \mathbb{N}$，由Abel变换知对$\forall m\geqslant k$，有
\begin{align*}
\sum\limits_{n=k+1}^{m} a_n = \sum\limits_{n=k+1}^{m} \frac{n a_n}{n} = \sum\limits_{n=k+1}^{m-1} \left( \frac{1}{n} - \frac{1}{n+1} \right) \sum\limits_{j=k+1}^{n} j a_j + \frac{1}{m} \sum\limits_{j=k+1}^{m} j a_j.
\end{align*}
令$m\to\infty$，结合$\sum\limits_{n=1}^{\infty} n a_n$收敛知
\begin{align}
\sum\limits_{n=k+1}^{\infty} a_n = \sum\limits_{n=k+1}^{\infty} \left( \frac{1}{n} - \frac{1}{n+1} \right) \sum\limits_{j=k+1}^{n} j a_j,\quad \forall k\in \mathbb{N}.
\label{eq:FHUAWAAWIWI2842}
\end{align}
并且对$\forall \varepsilon >0$，存在$K\in \mathbb{N}$，使得对
\begin{align}
\left| \sum\limits_{n=k+1}^{\infty} a_n \right| < \varepsilon,\quad \forall k\geqslant K.
\label{eq:FHUAWAAWIWI2843}
\end{align}
于是再结合\eqref{eq:FHUAWAAWIWI2842}式知，对$\forall k\geqslant K$，有
\begin{align}
\left| k \sum\limits_{n=k+1}^{\infty} a_n \right|
&\leqslant k \sum\limits_{n=k+1}^{\infty} \left( \frac{1}{n} - \frac{1}{n+1} \right) \left| \sum\limits_{j=k+1}^{n} j a_j \right| \nonumber\\
&= k \sum\limits_{n=k+1}^{\infty} \left( \frac{1}{n} - \frac{1}{n+1} \right) \left| \sum\limits_{j=k+1}^{\infty} j a_j - \sum\limits_{j=n+1}^{\infty} j a_j \right| \nonumber\\
&\leqslant k \sum\limits_{n=k+1}^{\infty} \left( \frac{1}{n} - \frac{1}{n+1} \right) \left( \left| \sum\limits_{j=k+1}^{\infty} j a_j \right| + \left| \sum\limits_{j=n+1}^{\infty} j a_j \right| \right) \nonumber\\
&\leqslant 2\varepsilon k \sum\limits_{n=k+1}^{\infty} \left( \frac{1}{n} - \frac{1}{n+1} \right) \leqslant \frac{2\varepsilon k}{k+1} < 2\varepsilon.
\label{eq:FHUAWAAWIWI2844}
\end{align}
从而由\eqref{eq:FHUAWAAWIWI2843}、\eqref{eq:FHUAWAAWIWI2844}式知，对$\forall k\geqslant K$，有
\begin{align*}
\left| \sum\limits_{n=1}^{\infty} n a_{n+k} \right| = \left| \sum\limits_{n=k+1}^{\infty} (n-k) a_n \right| \leqslant \left| \sum\limits_{n=k+1}^{\infty} n a_n \right| + \left| k \sum\limits_{n=k+1}^{\infty} a_n \right| \leqslant 3\varepsilon.
\end{align*}
故$\lim\limits_{k\to\infty} \sum\limits_{n=1}^{\infty} n a_{n+k} = 0$.
\end{enumerate}
\end{proof}




























\end{document}